\section{Appendix B. Implementation Details}

\subsection{B.1 Core Forward}

\begin{lstlisting}
def forward(self, hidden\_states, attention\_mask=None, **kwargs):
    B, N, D = hidden\_states.shape
    H, d = self.num\_heads, self.head\_dim
    
    # 1. Project
    q = self.q\_proj(hidden\_states).view(B, N, H, d).transpose(1, 2)
    k = self.k\_proj(hidden\_states).view(B, N, self.num\_kv\_heads, d).transpose(1, 2)
    v = self.v\_proj(hidden\_states).view(B, N, self.num\_kv\_heads, d).transpose(1, 2)
    
    # 2. RoPE
    if position\_embeddings is not None:
        q, k = apply\_rotary\_pos\_emb(q, k, cos, sin)
    
    # 3. GQA
    if self.kv\_repeat \textgreater{} 1:
        k = repeat\_kv(k, self.kv\_repeat)
        v = repeat\_kv(v, self.kv\_repeat)
    
    # 4. Adaptive E, F slicing
    L = N
    k\_eff = min(self.k\_proj\_dim, L)
    
    if L \textless{}= self.k\_proj\_dim:
        E\_act = self.E[:L, :L]
        F\_act = self.F[:L, :L]
    elif L \textless{}= self.max\_seq\_len:
        E\_act = self.E[:k\_eff, :L]
        F\_act = self.F[:k\_eff, :L]
    else:
        E\_act = F.interpolate(self.E[:k\_eff].unsqueeze(0), size=L, mode="linear").squeeze(0)
        F\_act = F.interpolate(self.F[:k\_eff].unsqueeze(0), size=L, mode="linear").squeeze(0)
    
    K\_proj = torch.matmul(E\_act, k)
    V\_proj = torch.matmul(F\_act, v)
    
    # 5. Attention
    scores = torch.matmul(q.float(), K\_proj.transpose(-2, -1).float()) * scale
    
    # 6. Causal fallback
    if self.is\_decoder and L \textgreater{} self.k\_proj\_dim:
        causal\_scores = torch.matmul(q.float(), k.transpose(-2, -1).float()) * scale
        causal\_mask = torch.tril(torch.ones(L, L, device=device, dtype=torch.bool))
        scores = causal\_scores.masked\_fill(~causal\_mask.unsqueeze(0).unsqueeze(0), -1e4)
        V\_proj = v
    
    # 7. Masking
    if attention\_mask is not None:
        scores = apply\_attention\_mask(scores, attention\_mask)
    
    # 8. Softmax + output
    attn = F.softmax(scores, dim=-1)
    context = torch.matmul(attn, V\_proj)
    context = context.transpose(1, 2).contiguous().view(B, L, self.all\_head\_size)
    
    if self.is\_nested\_encoder:
        out = context * self.output\_scale
    else:
        out = self.out\_proj(context * self.output\_scale)
    
    return out, None
\end{lstlisting}

\subsection{B.2 E, F Initialization}

\begin{lstlisting}
self.E = nn.Parameter(torch.empty(self.k\_proj\_dim, self.max\_seq\_len))
self.F = nn.Parameter(torch.empty(self.k\_proj\_dim, self.max\_seq\_len))

with torch.no\_grad():
    nn.init.normal\_(self.E, std=0.02)
    nn.init.normal\_(self.F, std=0.02)
    min\_dim = min(self.k\_proj\_dim, self.max\_seq\_len)
    self.E.data[:min\_dim, :min\_dim].fill\_diagonal\_(1.0)
    self.F.data[:min\_dim, :min\_dim].fill\_diagonal\_(1.0)
\end{lstlisting}

\subsection{B.3 Adaptive Slicing}

\begin{lstlisting}
if L \textless{}= self.k\_proj\_dim:
    # Short: identity-like
    E\_act = self.E[:L, :L]
elif L \textless{}= self.max\_seq\_len:
    # Standard Linformer
    E\_act = self.E[:k\_eff, :L]
else:
    # Extrapolation
    E\_act = F.interpolate(self.E[:k\_eff].unsqueeze(0), size=L, mode="linear").squeeze(0)
\end{lstlisting}

\subsection{B.4 Common Pitfalls}

\subsubsection{B.4.1 Forgetting RoPE Order}

RoPE must be applied \textbf{before} Linformer projection. After projection, 
position information is compressed and RoPE is meaningless.

\subsubsection{B.4.2 Breaking Causality}

Without the fallback, GPT-2 would see future tokens. Always check 
\texttt{is\_decoder and L \textgreater{} k}.

\subsubsection{B.4.3 Naive E, F Init}

Without diagonal init, $E K \approx 0$ and the model starts broken.

\subsubsection{B.4.4 Wrong Conv1D Orientation}

GPT-2's Conv1D stores (in, out). Transpose before chunking.

\subsubsection{B.4.5 Missing output\_scale}

\texttt{output\_scale} compensates for projection scaling. Initialize to 1.0.
\begin{lstlisting}

---
\end{lstlisting}